\begin{qfmethod}[Commodity swaps and index construction: using swaps to obtain or modify commodity exposure, and how weighting, rebalancing and roll rules affect index returns]\label{meth:commodity-swaps-index-construction}
\qffield{Statement} \emph{Commodity swaps.} A commodity swap is ``in essence a series of forward contracts on a commodity with different maturity dates and the same delivery prices'' \citep[p.~195]{s7-hulloptions}. A consumer of 100,000 barrels a year can pay \$5 million a year and receive 100,000 times the market price, locking in \$50 a barrel; a producer takes the opposite side (p.~781). Investors use swaps to:
\begin{itemize}
\item \emph{obtain exposure}: a \emph{total return swap} on a commodity index pays the index's total return in exchange for a floating rate plus a spread, giving index exposure without trading futures (Interpretation, from the index definitions in \hyperref[def:commodity-indices-risk-return]{index record});
\item \emph{modify exposure}: \emph{excess return swaps} pay the index excess return without the collateral leg; fixed-for-floating price swaps hedge a producer's or consumer's price; and \emph{basis swaps} exchange one price for another, such as a grade or location spread (Interpretation). \emph{Volatility swaps} exchange a preagreed volatility for realised volatility (Hull, p.~195).
\end{itemize}
Swap values follow from forward prices: the fixed payer's value is the discounted sum of forward price less fixed price times quantity.

\emph{Index construction} affects index returns through:
\begin{itemize}
\item \emph{weights}: production weights (S\&P GSCI, heavy in energy) or liquidity-and-production weights with caps (BCOM's 33 percent group cap), so sector returns dominate differently (\hyperref[def:commodity-indices-risk-return]{index record});
\item \emph{rebalancing}: fixed-weight rebalancing sells winners and buys losers, earning a \emph{diversification return} of roughly half the difference between average variance and average covariance \citep[pp.~24--25]{wide-erb-harvey-commodity}; Gorton and Rouwenhorst's monthly-rebalanced equal-weighted index embeds this strategy \citep[p.~9]{s16-r34};
\item \emph{roll rules}: which contracts are held and when they are rolled (fifth to ninth business day for the GSCI; sixth to tenth for BCOM) set the roll return; deferred or optimised contracts reduce roll costs in contango;
\item \emph{return variant}: spot, excess-return and total-return versions differ by the roll and the collateral yield.
\end{itemize}

\qffield{Practice} (Original; checked in the script.)
\begin{itemize}
\item \emph{Swap value}: a three-year swap at \$50 for 100,000 barrels a year, with forwards at 52, 54 and 55 and a 4 percent discount rate, is worth \$1.01 million to the fixed payer.
\item \emph{Total return swap}: an index total return of 7.1 percent against bills plus 25 basis points (4.25 percent) nets 2.85 percent.
\item \emph{Rebalancing}: one commodity alternates $+50$ and $-33.3$ percent (no net change), another is flat. Buy-and-hold gains nothing over four periods; an equal-weighted index rebalanced each period gains 8.51 percent, a 90/10 index 3.02 percent.
\item \emph{Roll rule}: in a contango costing 2 percent a month at the front but 0.75 percent further out, rolling the front contract loses 21.53 percent a year against 8.64 percent for deferred contracts.
\item \emph{Published check}: Erb and Harvey's heating oil and S\&P 500 example (1993--2003) shows a 10.9 percent rebalanced return against a 7.5 percent average, about 3.4 percent; half the difference between average variance (12.8) and covariance (5.3) is 3.75 percent, which they round to ``about 3.5\%'' (pp.~24--25).
\end{itemize}

\qffield{Conditions} The swap and rebalancing figures are illustrative; index rules are those in the 2026 methodologies.

\qffield{Pitfalls} Comparing indices without matching the return variant. Treating the diversification return as a risk premium (Erb and Harvey, p.~25). Ignoring counterparty risk in swaps.

\qffield{See also} \hyperref[def:commodity-indices-risk-return]{index record}; \hyperref[meth:commodity-futures-return-components]{return components}; \hyperref[mech:etf-creation-redemption]{ETF creation}.

\end{qfmethod}

\begin{qfdefinition}[Commodity indices and their record: the CRB, S\&P GSCI and Bloomberg indices, commodity ETFs and their energy weight, and the 1970--2006 risk, return and correlation figures Gliner reproduces]\label{def:commodity-indices-risk-return}
\qffield{Statement} The Commodity Research Bureau (CRB) Index ``is composed of 22 commodities and dates back to the 1950s''; the Deutsche Bank ETFs DBA (agriculture: corn, soybeans, sugar, cattle, coffee, cocoa and more) and DBC (energy, precious metals, agriculture and industrial metals) give basket exposure, the caution with DBC being that ``more than 50 percent of the index is energy weighted'' \citep[chap.~10, sec.\ ``CRB Index'', PDF p.~299]{s7-gliner}. ``Commodities have a higher volatility than foreign exchange, equities, and fixed income'', with ``an attractive Sharpe ratio as a relative comparison'' and, before 2008, negative correlation with equities; Table~10.1, from Goldman Sachs's \emph{Commodities as an Asset Class} (July 2007), gives annualised monthly continuously compounded returns, volatilities and correlations of equities, bonds, cash and the S\&P GSCI for 1970--2006, with an equity--GSCI correlation of $-0.04$ (sec.\ ``Return and Volatility'', p.~300).

\qffield{Notation} A production-weighted index (S\&P GSCI) weights commodities by world production value; the Bloomberg Commodity Index (BCOM, formerly DJ-UBS) uses liquidity and production with caps.

\emph{BCOM}: target weights combine two-thirds liquidity and one-third production percentages, using five-year averages, then apply annual diversification limits (no commodity above 15 percent, no commodity with its derivatives above 25 percent, no group such as energy above 33 percent, and at least 2 percent per commodity ``as liquidity allows''); the 2026 target weights are energy 29.44, grains 21.15, precious metals 18.84, industrial metals 15.76, softs 9.17 and livestock 5.64 percent, across 23 commodities; BCOM itself is an excess-return index, BCOMTR adds the return on Treasury-bill collateral (the 13-week bill auction high rate) and BCOMSP is a spot version; the roll shifts 20 percent a day over the sixth to tenth business days of the month; values before the launch on 14 July 1998 are ``historical estimations'' \citep[secs.~1.2 with Table~1, 2.2, 2.6, 3.1--3.2 and app.~B, pp.~4--6, 11, 23--25, 31--33 and~40]{bcom-methodology-2026}.

\emph{S\&P GSCI}: weights are contract production weights built from five-year world production averages, subject to total-dollar-value trading and reference-percentage-dollar-weight thresholds (at least 1.00 percent to enter, 0.10 percent to stay); the excess-return index (S\&P GSCI ER) ``is \dots\ calculated based on the contract daily return'' and the total-return index ``reflects the performance of a total return investment in commodities---contract daily return plus the daily interest on the funds hypothetically committed to the investment''; the roll runs from the fifth to the ninth business day of the month at 20 percent a day; the headline index has base date 31 December 1969 (value 100) and launch date 11 April 1991 \citep[PDF pp.~7, 9, 13 and~18, printed pp.~6, 8, 12 and~17]{spdji-gsci-methodology-2026}.

\qffield{Conditions} Figures to 2006; ETF details as of 2013. The 1970--2006 statistics are one sample from a bank's 2007 note. \emph{Live and back-tested history}. S\&P's education note states: ``The S\&P GSCI was launched April 11, 1991. \dots\ All data prior to such date is back-tested hypothetical data'' \citep[PDF p.~6, notes to the exhibits]{spdji-why-gsci-2024}, and the methodology's disclosure says that information before a launch date ``is hypothetical back-tested performance, not actual performance, and is based on the index methodology in effect on the launch date'' \citep[PDF p.~39]{spdji-gsci-methodology-2026}. The investable index ``arguably began in 1991'' and later indices were launched ``with backfilled history'' \citep[p.~36]{dunsby-nelson-2010}.

So Gliner's Table~10.1 (1970--2006) combines back-tested values from January 1970 to 10 April 1991 with live values from 11 April 1991; the boundary falls inside calendar 1991, which is neither wholly back-tested nor wholly live, and the monthly return for April 1991 straddles it. \emph{Variants.} A \emph{spot} index tracks futures prices only; an \emph{excess-return} index adds the roll from rolling futures; a \emph{total-return} index adds interest on fully collateralising the futures (S\&P's ER and TR definitions above; the education note separates the ``S\&P GSCI'' and the ``S\&P GSCI Excess Return (ER)'' by the roll, PDF p.~9).

Gliner's table labels the returns ``Annualized Monthly Returns (Continuously Compounded)'' and does not say which GSCI variant Goldman Sachs used, so its 1970--2006 figures should not be read as either total-return or spot returns. Whether correlations with risk assets rose significantly after 2008 is the financialisation question and is untested here.

\qffield{Practice} Read an index's rulebook for weights, roll schedule and rebalancing before using it as a benchmark; ETF returns differ from spot prices by the roll and fees (\hyperref[def:contango-backwardation-stocks]{contango and backwardation}).

\qffield{Pitfalls} Gliner's caution about energy exposure above 50 percent concerns DBC; the GSCI's energy weight has declined since 2008 (education note, Exhibit~8, PDF p.~8), while BCOM's 33 percent group cap holds its 2026 energy target at 29.44 percent. Whether commodities are particularly useful as inflation hedges is an empirical question, to be measured by horizon and regime.

\qffield{See also} \hyperref[def:contango-backwardation-stocks]{contango and backwardation}; \hyperref[def:macro-regime-asset-map]{regime map}; \hyperref[def:crude-oil-markets]{crude oil}.


\end{qfdefinition}
